\chapter{Introduction}

%\chapter*{Introducere}
\label{intro}

\section{Introducing the problem}

\par Public institutions play an important role in the daily lives of citizens, providing access to services such as official documents, permits, and administrative procedures. Because of this, it is important that the processes involved are efficient, reliable, and easy to use.

\par In many cases, these processes are still based on physical documents and in-person interactions. Citizens are often required to visit institutions, wait in queues, and submit printed files in order to complete simple requests. This approach leads to long waiting times, repeated visits, and inefficient handling of information.

\par In addition, physical documents can be lost, damaged, or incorrectly processed, which can delay important procedures. These inefficiencies do not only create inconvenience, but can also impact the ability of citizens to complete time-sensitive administrative tasks.

\par As a result, improving the way public institutions handle document-based workflows becomes an important problem to address.

\section{Role of Information Technology in Public Administration}

\par The use of information technology in public institutions has become increasingly important in recent years. Digital systems allow institutions to manage large amounts of data, automate processes, and reduce the time required to complete administrative tasks.

\par By moving services to digital platforms, interactions between citizens and institutions can be simplified. Instead of relying on physical documents and in-person visits, many processes can be handled online, improving accessibility and reducing waiting times.

\par In addition, digital systems make it easier to store, organize, and retrieve information. This can improve decision-making and reduce errors caused by manual processing of documents.

\par However, the adoption of such systems is not always straightforward. Public institutions may face challenges such as limited financial resources, lack of technical infrastructure, and insufficient digital skills among both employees and citizens.

\par Despite these challenges, the use of information technology remains a key factor in improving the efficiency and quality of public services.

\section{Challenges in Traditional Systems}

\par Traditional document handling in public institutions comes with several practical challenges that affect both citizens and employees.

\par One of the main issues is the need for physical presence. In most cases, citizens are required to visit an institution in person in order to submit or collect documents. This creates queues and waiting times, especially during busy hours or peak periods.

\par Another problem is the use of paper-based documents. Files are printed, copied, and stored physically. This takes space and also creates the risk of documents being lost, damaged, or misplaced. It also makes it harder to search and manage information when needed.

\par The process of checking and validating documents is often manual. Employees need to go through each file one by one, which takes time and can lead to delays, especially when the number of requests is high.

\par Feedback to citizens is also not always immediate. If a document is missing or incorrect, the citizen usually finds out later, sometimes after waiting for a long time or making another visit. This leads to repeated submissions and additional waiting time.

\par In addition, many of these processes are not transparent to the user. Citizens often do not know the current status of their request or how long it will take until it is completed.

\par These challenges show that the traditional way of handling document-based workflows is slow and can be improved with more structured and digital solutions.

\section{Limitations of existing digital solutions}

\par Even though some public services already use digital platforms, they are often limited in scope. Most of these systems are built for a specific type of service or a single institution.

\par In many cases, users still need to switch between different systems depending on the service they need. This creates fragmentation and makes the overall process less consistent.

\par Some systems also only cover part of the workflow, such as submitting documents or signing them, but not the full process from submission to final approval and feedback.

\par Because of this, there is still a need for systems that can be reused across different institutions, where each institution can use its own instance of the same workflow structure instead of relying on completely separate tools.

\section{Proposed solution}

\par The proposed system, Doclane, is a web-based application designed to handle document-based workflows between citizens and public institutions.

\par The system allows users to submit requests, upload required documents, and track the status of their applications. On the other side, institution employees can review submissions, validate documents, and provide feedback when needed.

\par The goal of the system is to reduce the need for physical visits and repeated submissions by moving the process into a structured digital workflow.

\par Doclane also focuses on making the process more clear for users by showing what is required at each step and giving feedback when something is missing or incorrect.

\section{Objectives}

\par The main objective of this thesis is to design and implement a system that improves document-based workflows in public institutions.

\par The specific objectives include:
\begin{itemize}
    \item Designing a backend system that can handle document requests.
    \item Implementing secure authentication and role-based access.
    \item Building a structured workflow for submission and validation.
    \item Using cloud storage for document handling.
    \item Making the system scalable and easy to extend.
\end{itemize}

\section{Scope of the system}

\par The system focuses on improving the workflow between citizens and public institutions when it comes to document-based requests.

\par It does not change the legal process of document approval. Final decisions are still made by authorized personnel within the institution.

\par The system is designed as a support tool that helps organize and digitalize the existing workflow, not replace institutional rules or procedures.

% \section{Motivation}

% \par The idea behind Doclane was optimizing a complex, lengthy process that most people go through on a regular basis into a predominantly digital, simplified and shorter one.

% \par With the way public institutions are structured, they handle thousands of requests per day, many of those requiring files that will be kept in a physical archive for years, or maybe indefinetly. This is not only wasteful when taking into account the consumed paper, but also in human time.

% \par Paper can be lost and it can be damaged. The solution to making sure it won't get lost is to make copies of it, which is again wasteful. Why not optimize this process by digitalizing the biggest bottleneck in it and implement modern technologies such as cloud storage, which proves to be so much more efficient in the long term? And why stop at just cloud storage, when we can optimize this process, to also save time instead of physical resources? We reduce costs, save the planet and save our time, everyone wins in this equation.

% \section{Problem statement}

% \par This thesis addresses inefficiencies in traditional public administration workflows, particularly in the submission, validation, and processing of document-based requests. These workflows are often characterized by manual handling, repeated physical interactions, and lack of structured feedback mechanisms.

% \section{Objectives}

% The main objective \cite{Sommerville2010} of this thesis is to design and implement a cloud-based system that improves the efficiency of document request workflows in public institutions.

% The specific objectives include:
% \begin{itemize}
%     \item Designing a scalable backend architecture for handling document requests.
%     \item Implementing secure authentication and role-based access control.
%     \item Providing a structured document submission and validation workflow.
%     \item Integrating cloud storage for secure file management.
%     \item Ensuring system reliability through modern cloud infrastructure.
% \end{itemize}

% \section{Limitations}

% The system is designed as a digital workflow optimization platform and does not replace the legal authority of public institutions. Final document issuance and validation remain under institutional control.

% \section{Current process}

% \par Let's take the example of how you could, in this current day, retrieve a document from a public institution, maybe an urbanim certificate so that you can build stuff on your property. You'd schedule an appointment whenever the next free time slot is (and queue times are definetly not short), then you go present yourself, you take your identification with you and all the required documents such that it can be emitted. Let's say you did not forget anything and you do not need to reschedule, although this happens sometimes too. By the time you have your certificate emitted, half your day is gone, not only because the process itself is lengthy but due to the sheer number of other citizens that need their requests solved you might have to wait longer than your appointment timeline, and maybe delays occur too.

% \section{What Doclane promises}

% \par Doclane digitalizes the bottleneck in this entire process. Let's take a look at how you can receive the same urbanism certificate while reducing human interaction to just one step. You wake up, you log in to Doclane, your account is linked to the local institution you belong to, and in the app you select the certificate you need, and then choose to submit the request. You will be prompted by the application to upload relevant documents to your applications in specific slots, each slot also providing you with examples on how the uploads should look. You uploaded everything within 5 minutes, then on the other side public institution employees are notified, they check the documents, emit an approval and you go pick up your document.

% \par Did you submit something that the institution cannot accept? They can notify you within the application and let you know what you didn't do right, so that you may correct it. No more re-scheduling appointments and losing hours upon hours to have the paper you need released.